\section{Por qué la evaluación de agentes de próxima generación no puede basarse solo en la precisión}

Los benchmarks tradicionales comprimen la tarea en una entrada estática y una respuesta corta, adecuados para medir conocimiento y razonamiento local, pero difícilmente reflejan la capacidad del agent de planificar, invocar herramientas, recuperarse de errores y entregar productos complejos a lo largo de decenas de minutos u horas. Esta clase responde tres preguntas mediante tres evaluaciones complementarias:

\begin{enumerate}
    \item \textbf{METR Time Horizon}: ¿cuánto tiempo de una tarea humana puede completar el modelo de forma autónoma con una confiabilidad dada?
    \item \textbf{GDPval}: ¿tiene la salida del modelo, comparada con la de expertos reales de la industria, un valor económico de uso?
    \item \textbf{DeepScholar-Bench}: ante una tarea de investigación abierta, ¿puede el sistema recuperar información de manera exhaustiva, organizar el conocimiento y ofrecer citas verificables?
\end{enumerate}

\videofigure{figures/measuring-ai-progress.jpg}{Evaluar el progreso de la IA requiere considerar simultáneamente la complejidad de la tarea, la confiabilidad y el valor económico; los benchmarks tradicionales se saturan rápidamente.}{00:02:54--00:04:26}

\begin{importantbox}{La capacidad del agent tiene al menos tres ejes independientes}
La duración de la tarea responde «cuánto tiempo puede hacer», la probabilidad de éxito responde «qué tan confiable es» y la calidad de la salida responde «si vale la pena usarlo». Un solo puntaje no puede expresar estos tres ejes a la vez.
\end{importantbox}

\begin{warningbox}{El puntaje de un benchmark no es la proporción real de automatización}
Las tareas experimentales suelen ser más limpias, con un contexto más completo y herramientas más estables, y rara vez incluyen comunicación organizacional, conocimiento tácito y restricciones de responsabilidad. Que un modelo tenga éxito en un benchmark no significa que la misma proporción de puestos reales pueda completarse sin supervisión humana.
\end{warningbox}

\subsection{Resumen de la sección}

Lo que se evalúa en los agents de largo plazo ya no es una respuesta, sino una trayectoria de ejecución y un producto de trabajo final. El tiempo, la confiabilidad, el valor económico, la dependencia del contexto y la verificabilidad deben integrarse juntos en el sistema de métricas.

